\section{Scope of inference}
\label{sec:limitations}

\paragraph{Model and evidence join.}
The trace threshold is a swept candidate rather than the measured crossover of
the resident-policy source, and the mechanism's $H=32$ horizon is absent from
the trace model. The two result sets therefore cannot be multiplied. The
offline optimum also assumes future knowledge, zero service time, unlimited
capacity, no upper batch size, and deadlines on launch rather than completion.
The safe-suffix abstraction extrapolates beyond tested cohort sizes and does
not cover non-monotone benefit, arbitrary deadlines, or batch-size-dependent
service.

\paragraph{Trace scope and executable identity.}
The replay conditions on one \TraceSessions-session panel from three related
customer-service domains and imposes stationary Poisson arrivals. Its three
seeds measure Monte Carlo variation under that panel and model, not a workload
population. Bursts, correlated releases, and other corpora can move the
boundary. The outcome-derived route key omits state-machine node, schema,
arguments, policy context, and multi-tool identities, so it is a conditioning
proxy rather than verified semantic fusion. The packing arrays also omit a
per-session sequence constraint; overlapping spans and completion-order
inversions can affect the wider deadline surface.

\paragraph{Mechanism and sampling scope.}
Resident-policy-001 makes one global binary decision over a regular synthetic
state array. A route service still needs per-event compaction, variable route
bodies, ingress and egress, effect ordering, recovery, and CPU fallback. Setup,
state reset, result copy, and validation are outside timing. The four named
placements comprise two L4s, one H100, and one GTX; GPU, provider, host, image,
driver, and region are confounded. The local full run reuses the development
GPU, and one external placement extends the frozen first-stage scope. Timing
rows are technical repetitions, so the current performance result supports
named-placement effect sizes rather than a hardware-population inference.

\paragraph{Deployment and algorithmic scope.}
The mechanism metric is batch-average cohort-horizon wall time, not invocation
P99 or end-to-end task time. A tuned CPU implementation of the real transition,
CPU core-seconds, energy, cost, model throughput, TTFT, TPOT, task utility, and
external-effect reliability remain deployment measurements. The dynamic
program is a specialized exact instrument within mature batching, clustering,
and scheduling literatures; the paper makes no exhaustive algorithmic priority
claim.
